%-------------------------------------------------------------------------------
% Cover letter — Quantitative research summer internships
% Zeyu (Steven) Zhang.  Compile with: tectonic cover-letter-quant.tex
%
% Before sending, replace the three [BRACKETED] placeholders:
%   1. [Company] and the recipient line
%   2. [ROLE / PROGRAM NAME] in the opening paragraph
%   3. The company-specific sentence in the closing paragraph
%-------------------------------------------------------------------------------

\documentclass[10.5pt]{article}

\usepackage[letterpaper, margin=2.2cm, top=1.8cm]{geometry}
\usepackage[T1]{fontenc}
\usepackage[utf8]{inputenc}
\usepackage{XCharter}
\usepackage[dvipsnames]{xcolor}
\usepackage{fontawesome5}
\usepackage{hyperref}

\definecolor{accent}{HTML}{1F3A6E}
\definecolor{bodygray}{HTML}{2B2B2B}
\definecolor{metagray}{HTML}{5A5A5A}
\colorlet{rulegray}{black!25}

\hypersetup{colorlinks=true, urlcolor=accent, linkcolor=accent,
  pdftitle={Zeyu Zhang - Cover Letter (Quantitative Research)}, pdfauthor={Zeyu Zhang}}

\color{bodygray}
\pagestyle{empty}
\setlength{\parindent}{0pt}
\setlength{\parskip}{9pt}

\begin{document}

%--------------------------- Header --------------------------------------------
\begin{center}
  {\Huge\scshape Zeyu (Steven) Zhang}\\[6pt]
  {\color{metagray}\itshape Ph.D.\ Student in Statistics and Data Science, Northwestern University}\\[6pt]
  {\small
    \faEnvelope\ \href{mailto:zeyuzhang2028@u.northwestern.edu}{zeyuzhang2028@u.northwestern.edu}
    \;$|$\;
    \faGlobe\ \href{https://zeyuzhang1901.github.io}{zeyuzhang1901.github.io}
    \;$|$\;
    \faLinkedin\ \href{https://www.linkedin.com/in/zeyu-zhang-42813a271/}{LinkedIn}
    \;$|$\;
    \faGraduationCap\ \href{https://scholar.google.com/citations?user=dr-xetEAAAAJ}{Google Scholar}
  }\\[4pt]
  {\color{accent}\rule[0.6pt]{2.4cm}{1.6pt}}{\color{rulegray}\rule[1.15pt]{\dimexpr\textwidth-2.4cm\relax}{0.5pt}}
\end{center}

\vspace{4pt}
{\color{metagray}\today}

Dear [Company] Quantitative Research Recruiting Team,

I am a third-year Ph.D.\ student in Statistics and Data Science at Northwestern
University, applying for the [ROLE / PROGRAM NAME]. My research is about a question
every systematic strategy must answer before capital is at risk: \textbf{when can a
backtest be trusted?} I study look-ahead bias in evaluations of LLM-based
forecasters---models that have already ``read the future'' in their training
data---and build the statistical machinery to detect, measure, and remove it.

My most recent paper (in submission to TMLR) shows that the standard defense,
comparing scores before and after a training cutoff, is uninformative: legitimate
recency reproduces the leakage signature, and I prove that no passive backtest can
separate contamination from genuine skill. One defensible external reference
restores measurement---difference-in-differences against a matched clean control,
or regression discontinuity at a known cutoff---yielding a leakage-adjusted score
with confidence intervals. I did not stop at theory: I validated the estimators by
planting leakage at controlled doses in twin models, where they recovered the
injected dose and returned null on the clean twin, then audited frontier models on
1{,}646 prediction-market questions, clearing five whose apparent forecasting edge
was recency alone. The habits this work demands---distrusting in-sample signals,
designing falsification checks, pre-registering analyses---are the working habits
of quantitative research.

The rest of my record shows I can carry a problem from proof to large-scale
experiment. My first-author NeurIPS 2023 paper cast learning-to-rank from logged
clicks as offline reinforcement learning: counterfactual learning from biased
logged feedback, structurally the same problem as learning from historical trading
data. Two companion papers under review extend the backtesting program to
mitigation---constrained RL post-training with convergence guarantees (NeurIPS
2026) and an inference-time auditing system evaluated on 2{,}769 forecasting
instances (EMNLP 2026). I work daily in Python and PyTorch, build evaluation
pipelines end to end, and have taught six quarters of statistics and data science
as a TA, so I am comfortable explaining technical decisions to a non-specialist
audience.

[One sentence on why this firm specifically: a team, a research area, or an aspect
of its culture that matches the work above.] I would welcome the chance to bring
this skepticism and rigor to problems where the data are noisier and the stakes are
real. Thank you for your consideration.

\vspace{6pt}
Sincerely,\\[2pt]
{\scshape Zeyu (Steven) Zhang}

\end{document}
